\chapter{HARDWARE REFERENCE}
\label{ch:platform}

This chapter describes every piece of hardware on the robot: what it is, its
relevant specifications, and how it is connected. The next chapter covers
the software changes that were needed to make each piece work on the
BeagleBone AI.

\section{CHASSIS AND DRIVETRAIN}

The chassis is a small four-wheel-drive platform bought from SparkFun, with an
aluminium plate added as a deck for the electronics. Each wheel is driven by
its own brushed DC gear motor, and the two motors on each side are commanded
together, which makes the robot a skid-steer vehicle
(\S\ref{sec:skid}). Table~\ref{tab:chassis} gives the measured geometry.

\begin{table}[htbp]
\centering
\caption{Chassis geometry (measured 2026-10-01)}
\label{tab:chassis}
\small
\begin{tabular}{lll}
\toprule
Quantity & Value & Note \\
\midrule
Wheel diameter & \SI{79.2}{mm} & radius $r=\SI{39.6}{mm}$ \\
Track (left--right wheel centres) & \SI{120}{mm} & $b$ \\
Wheelbase (front--back axles) & \SI{101}{mm} & $L$; $L/b=0.84$ \\
Encoder resolution & 15 counts per wheel turn & \SI{16.6}{mm} of travel per count \\
Top speed (wheels lifted, duty 1) & about \SI{1.2}{m/s} & capped to about \SI{0.65}{m/s} in software \\
\bottomrule
\end{tabular}
\end{table}

The encoders are DAGU magnetic encoder kits: an 8-pole ring magnet on each
wheel hub read by an Allegro A1230 dual Hall sensor. Because the magnet turns
with the wheel, the encoder measures the wheel itself, after the gearbox. This
is the opposite trade from the TWIP, whose encoders were on the motor shaft
ahead of the gearbox (thesis \S4.2.3): here gear backlash is seen by the
encoder, but the resolution is about 100 times coarser per wheel turn.

\section{BEAGLEBONE AI}
\label{sec:bbai}

The BeagleBone AI is built around the Texas Instruments AM5729 system on chip
(the Sitara AM57x family). It keeps the BeagleBone Black's form factor and its
two 46-pin expansion headers, P8 and P9, so BeagleBone capes physically fit.
Table~\ref{tab:bbai} summarises the parts that matter for the robot.

\begin{table}[htbp]
\centering
\caption{BeagleBone AI}
\label{tab:bbai}
\small
\begin{tabularx}{\linewidth}{>{\raggedright\arraybackslash}p{1.45in}X}
\toprule
Block & Details and use on the robot \\
\midrule
CPU & 2$\times$ ARM Cortex-A15 at \SI{1.5}{GHz}. Runs Linux, ROS and all the
  Python nodes. \\
DSP & 2$\times$ TI C66x floating-point DSP. Unused. \\
EVE & 4$\times$ Embedded Vision Engine vector processors, used by TI's TIDL
  deep-learning library. Used for the 2022 webcam classification demo; the
  OAK-D-Lite has since taken over vision. \\
PRU-ICSS & 2 industrial communication subsystems, each with 2 PRU cores at
  \SI{200}{MHz} (4 in total). PRU1\_1 counts encoder 4; PRU2\_0 drives motor 1's
  direction pin (\S\ref{sec:pru}). \\
IPU & 2$\times$ dual Cortex-M4. Unused. \\
Memory & \SI{1}{GB} DDR3, \SI{16}{GB} eMMC (boots Debian), microSD slot. \\
Peripherals & eQEP quadrature decoders (3 used for encoders 1--3), ePWM
  (motor PWM), UARTs, I$^2$C, SPI, GPIO, ADC. \\
Connectivity & Gigabit Ethernet; AzureWave AW-CM256SM module (Broadcom BCM43455):
  dual-band Wi-Fi on SDIO and Bluetooth on UART6; USB Type-C (power and USB~3
  device/host); USB Type-A host (USB~2, used by the OAK-D-Lite). \\
Power & \SI{5}{V} in through USB-C, or through the headers' \texttt{SYS\_5V}
  pins, which is how the Robotics Cape powers it. A TPS659037 PMIC makes the
  internal rails. \\
Cooling & A fan (from the parts list) on the heat sink; the AM5729 runs hot
  under load. \\
\bottomrule
\end{tabularx}
\end{table}

\paragraph{Header multiplexing.} On the BeagleBone Black each header pin
connects to one ball of the AM335x. On the BeagleBone AI many header pins are
wired to \emph{two} AM5729 balls (written P9.19a and P9.19b, for example), to
reach more functions. Only one of the two may be driven; the other must be set
to an input or to mode 15 (off) in the device tree, or the two outputs fight.
Each ball has up to 16 multiplexer modes; the robot uses mode 10 for eQEP and
ePWM, mode 12 and 13 for PRU inputs and outputs, mode 14 for GPIO and modes
1--7 for UART and I$^2$C. The PRU pin map that guided the motor~1 fix is in
\file{claudefiles/beagleboneAI/PRUs/} in the project files: it shows that
P9.13's reachable ball (C17) offers \texttt{pr2\_pru0\_gpo15} in mode 13 and
no ordinary GPIO.

\section{ROBOTICS CAPE}
\label{sec:cape}

The Robotics Cape (Strawson Design SD-101D, revision D, sold as
BB-CAPE-ROBOTICS) was designed for the BeagleBone Black. It carries the
robot's power system, motor drivers, IMU and connectors.
Figure~\ref{fig:cape} shows its blocks, and the schematic is in
\file{claudefiles/RoboticsCape/} in the project files.

\begin{figure}[htbp]
\centering
\begin{tikzpicture}[font=\scriptsize,
  b/.style={draw,rounded corners=2pt,fill=white,align=center,minimum height=9mm,
            minimum width=25mm,inner sep=2pt},
  p/.style={b,fill=orange!10},
  s/.style={b,fill=blue!6},
  a/.style={-{Latex[length=1.6mm]},semithick}]
\node[p] (jack) at (0,3.4) {Barrel jack\\ \SIrange{9}{18}{V} \texttt{VDC}};
\node[p] (lipo) at (0,1.8) {2S LiPo, JST-XH\\ \texttt{V\_BATT} \SIrange{6}{8.4}{V}};
\node[p] (chg) at (3.4,3.4) {Charger MP2615\\ \SI{1}{A}, \SI{4.2}{V}/cell};
\node[p] (prot) at (3.4,1.8) {Balancer BQ29209\\ protector S-8261};
\node[p] (reg5) at (7.0,3.4) {\SI{5}{V} buck AP1509\\ \SI{2}{A}};
\node[p] (reg6) at (7.0,1.8) {\SI{6}{V} buck AOZ1284\\ servo rail};
\node[s] (bb) at (10.6,3.4) {BeagleBone AI\\ \texttt{SYS\_5V}, P8/P9};
\node[s] (servo) at (10.6,1.8) {8 servo/ESC\\ outputs (unused)};
\node[s] (hb) at (0,-0.2) {2$\times$ TB6612FNG\\ 4 H-bridges M1--M4};
\node[s] (enc) at (3.4,-0.2) {4 encoder ports\\ E1--E4 (JST-SH)};
\node[s] (imu) at (7.0,-0.2) {MPU-9250 IMU\\ BMP280 barometer};
\node[s] (io) at (10.6,-0.2) {UART1, UART5, GPS,\\ DSM2, SPI, I$^2$C, ADC};
\draw[a] (jack) -- (chg);
\draw[a] (chg) -- (prot);
\draw[a] (lipo) -- (prot);
\draw[a] (jack.north) -- ++(0,0.4) -| (reg5.north);
\draw[a] (prot.east) -- ++(0.35,0) |- (reg5.west);
\draw[a] (prot) -- (reg6);
\draw[a] (reg5) -- (bb);
\draw[a] (reg6) -- (servo);
\draw[a] (lipo.south) -- (hb.north) node[midway,left]{\texttt{V\_BATT}};
\end{tikzpicture}
\caption{Robotics Cape (SD-101D) blocks}
\label{fig:cape}
\end{figure}

\subsection{Power}

\begin{description}
\item[Inputs.] A 2-cell LiPo on a 3-pin JST-XH connector (\texttt{V\_BATT}
  plus the cell midpoint for balancing), and a \SI{5.5}{mm}$\times$\SI{2.1}{mm} centre-positive
  barrel jack for \SIrange{9}{18}{V} DC (\texttt{VDC}). A \SI{9}{V}, \SI{2}{A}
  adapter is on the parts list. Either input can run the robot.
\item[Charger.] An MPS MP2615 charges the pack from \texttt{VDC} at \SI{1}{A}
  (set by a \SI{0.1}{\ohm} sense resistor) to \SI{4.2}{V} per cell.
\item[Protection.] A TI BQ29209 balances the two cells and a Seiko S-8261
  cuts the pack off at \SI{4.3}{V} per cell (overcharge) or \SI{2.5}{V} per
  cell (over-discharge).
\item[\SI{5}{V} rail.] A Diodes AP1509 buck regulator, fed from \texttt{VDC} or
  \texttt{V\_BATT} through diodes, makes \SI{5}{V} at up to \SI{2}{A}. This
  rail powers the BeagleBone through \texttt{SYS\_5V}, and through it the
  BeagleBone's USB ports and everything plugged into them. Its \SI{2}{A} limit
  is central to the board resets (\S\ref{sec:powerbudget}).
\item[\SI{6}{V} rail.] An AOZ1284 buck powers the servo headers when enabled
  by \texttt{SERVO\_EN}. Unused.
\item[Monitoring.] Two dividers of \SI{47}{k\ohm} and \SI{4.7}{k\ohm} scale
  \texttt{VDC} and \texttt{V\_BATT} by $1/11$ into the \SI{1.8}{V} ADC.
  Four LEDs show battery level.
\end{description}

\subsection{Motor drivers}

Two Toshiba TB6612FNG dual H-bridges drive the four motor ports (H1: M1 and
M2; H2: M3 and M4) directly from \texttt{V\_BATT}. Each channel has two
direction inputs (\texttt{IN1}, \texttt{IN2}, the cape's \texttt{MDIRxA/B})
and a PWM input, and a shared standby line (\texttt{MOT\_STBY}) enables both
chips. Per channel the TB6612FNG is rated \SI{1.2}{A} continuous and
\SI{3.2}{A} peak, at up to \SI{15}{V}. Its truth table is
\begin{center}
\small
\begin{tabular}{cccl}
\toprule
\texttt{IN1} & \texttt{IN2} & \texttt{PWM} & Output \\
\midrule
H & L & H & forward (current one way) \\
L & H & H & reverse \\
H & H & x & short brake \\
L & L & H & coast (outputs off) \\
x & x & L & short brake \\
\bottomrule
\end{tabular}
\end{center}
so speed is set by PWM duty with the direction held on \texttt{IN1/IN2}
(sign--magnitude drive). During the PWM off-time the bridge short-brakes the
motor, which makes average speed close to proportional to duty.
librobotcontrol runs the PWM at \SI{25}{kHz}.

\subsection{Encoder ports}

Four 4-pin JST-SH connectors, E1--E4, each carry GND, \SI{3.3}{V} and the A
and B signals. E1--E3 go to three eQEP hardware decoders; E4 goes to PRU
inputs (on the BeagleBone Black it was read by PRU firmware too, since the
AM335x has only three eQEPs). There are no pull-up resistors on these lines on
the cape, which is why the device tree must enable the SoC's internal ones
(\S\ref{sec:dt}).

\subsection{Sensors on the cape}

The MPU-9250 IMU sits on the cape's I$^2$C2 bus with its interrupt on
\texttt{IMU\_INT} (P9.25). A Bosch BMP280 barometer shares the bus at address
\texttt{0x76}; it is not used. The I$^2$C buses have \SI{4.7}{k\ohm} pull-ups.

\subsection{Connectors}

\begin{table}[htbp]
\centering
\caption{Robotics Cape connectors and what is plugged into them}
\label{tab:capeconn}
\small
\begin{tabularx}{\linewidth}{lllX}
\toprule
Connector & Type & Pins & On this robot \\
\midrule
M1--M4 & JST-ZH 2-pin & motor $\pm$ & M1 left back, M2 left front, M3 right front, M4 right back \\
E1--E4 & JST-SH 4-pin & GND, 3.3V, A, B & encoders: ch1 left back, ch2 right back, ch3 left front, ch4 right front \\
UART1 & JST-SH 4-pin & GND, 3.3V, RX, TX & RPLIDAR A1 data (\texttt{/dev/ttyS3}) \\
GPS & JST-SH 6-pin & GND, 3.3V, RX, TX, 5V, GND & EM-506 (\texttt{/dev/ttyS2}) \\
UART5 & JST-SH 4-pin & GND, 3.3V, RX, TX & free \\
DSM2 & JST-ZH 3-pin & RX, GND, 3.3V & free (radio receiver was considered) \\
I$^2$C1, SPI1.1/1.2, ADC & JST-SH & & free \\
PWR & JST-SH 4-pin & 5V, 3.3V, GND & free; auxiliary power \\
Servo 1--8 & 0.1'' 3-pin & signal, 6V, GND & free \\
LiPo / jack & JST-XH 3-pin / barrel & & battery / bench supply \\
\bottomrule
\end{tabularx}
\end{table}

The cape's buttons (Pause, Mode) and LEDs (red, green, battery) are routed to
P8 GPIOs; the buttons do not work on the BeagleBone AI yet.

\section{SENSORS}

\subsection{MPU-9250 IMU}

\begin{table}[htbp]
\centering
\caption{MPU-9250 characteristics used here (datasheet rev.~1.1)}
\label{tab:mpu}
\small
\begin{tabularx}{\linewidth}{l>{\raggedright\arraybackslash}X>{\raggedright\arraybackslash}p{40mm}}
\toprule
& Value & Note \\
\midrule
Gyro range & \SI{\pm250}{\degree/s} (131 LSB per \si{\degree/s}) & default in librobotcontrol \\
Gyro rate-noise density & \SI{0.01}{\degree/s/\sqrt{Hz}} & \\
Gyro zero-rate offset & \SI{\pm5}{\degree/s} initial, \SI{\pm30}{\degree/s} over temperature & hence start-up bias removal \\
Accel noise density & \SI{300}{\micro g/\sqrt{Hz}} & \\
Magnetometer (AK8963) & \SI{\pm4800}{\micro T} range & axes differ from gyro; library realigns \\
\bottomrule
\end{tabularx}
\end{table}

\subsection{RPLIDAR A1M8}

A triangulation laser scanner on a rotating head driven by its own motor.
It measures up to 8000 times per second, \SIrange{0.15}{12}{m}, with
resolution under 1\% of distance, over 360 degrees. Scan rate is 1--10~Hz
(5.5~Hz typical); on this robot it settles at about \SI{7.2}{Hz}. The scanner
core needs \SIrange{4.9}{5.5}{V} with low ripple (about \SI{300}{mA} running,
\SI{500}{mA} at start-up) and the motor \SIrange{5}{9}{V}. Data is
\SI{115200}{baud} serial at \SI{3.3}{V}, which goes to the cape's UART1
header; the cape's UART headers carry only \SI{3.3}{V}, so the lidar's power
comes from a separate \SI{5}{V} feed.

\subsection{EM-506 GPS}

A GlobalSat module with a SiRF Star IV chipset and built-in patch antenna:
\SI{-163}{dBm} tracking sensitivity, under \SI{2.5}{m} horizontal error (50\%,
open sky), cold start under \SI{35}{s} (hot under \SI{1}{s}), NMEA 0183 at
\SI{4800}{baud}, \SI{1}{Hz}, SBAS (WAAS) capable. It draws about
\SI{50}{mA} from \SIrange{4.5}{6.5}{V}, which the cape's GPS header provides.

\subsection{Wheel encoders (A1230)}

Each A1230 contains two Hall switches \SI{1}{mm} apart whose outputs are in
quadrature when they face a matching ring magnet. Outputs are open-drain,
sinking up to \SI{20}{mA}, and need an external pull-up; supply is
\SIrange{3.3}{18}{V}. Decoding the four edges of the two channels gives
$4\times$ the number of pole pairs per revolution.

\subsection{OAK-D-Lite camera}

A stereo depth and AI camera; Chapter~\ref{ch:camera} covers it in detail.

\subsection{PS4 controller}

A Sony DualShock~4 over Bluetooth Classic HID, read through the kernel's
\texttt{hid-sony} driver as \texttt{/dev/input/js0}.

\section{POWER BUDGET}
\label{sec:powerbudget}

Everything except the motors runs from the cape's \SI{5}{V} rail, which is
limited to \SI{2}{A} by the AP1509. Table~\ref{tab:budget} estimates the
load.

\begin{table}[htbp]
\centering
\caption{\SI{5}{V} rail load (estimates from datasheets)}
\label{tab:budget}
\small
\begin{tabularx}{\linewidth}{>{\raggedright\arraybackslash}X l l}
\toprule
Load & Current at \SI{5}{V} & Source \\
\midrule
BeagleBone AI, both A15 cores busy, Wi-Fi and BT & \SIrange{0.8}{1.2}{A} & typical for the board \\
OAK-D-Lite running depth and detection & \SIrange{0.8}{1.2}{A} (1.5 max) & datasheet: \SIrange{4}{6}{W} \\
RPLIDAR A1 core + motor & \SIrange{0.4}{0.6}{A} & datasheet \\
GPS, IMU, encoders & \SI{0.07}{A} & datasheets \\
\midrule
Total & \SIrange{2.1}{3.1}{A} & vs.\ \SI{2}{A} regulator \\
\bottomrule
\end{tabularx}
\end{table}

The estimated total exceeds the regulator's rating once the camera runs
alongside everything else, which matches what was observed: the camera
alone, or everything without the camera, ran fine, while the full load reset
the board with the battery voltage steady. The camera datasheet anticipates
this: hosts that cannot supply its current should use a USB Y-adapter that
feeds the camera from a separate supply. That is the planned fix (a USB
splitter cable and battery bank), which moves \SIrange{0.8}{1.5}{A} off the
cape's rail.

The motors draw from \texttt{V\_BATT} directly through the H-bridges, not
from the \SI{5}{V} rail. Their start-up and stall currents can still sag the
battery and the regulator's input, which is another reason to keep the
\SI{5}{V} load well below \SI{2}{A}.
